% Part: first-order-logic
% Chapter: beyond
% Section: introduction

\documentclass[../../../include/open-logic-section]{subfiles}

\begin{document}

\olfileid{fol}{byd}{int}

\olsection{概述}

一阶逻辑并非唯一受关注的逻辑系统：它有许多扩展与变体。
一个逻辑系统通常包括语言的形式规范，通常（但不必然）包括一个推演系统，
以及通常（但不必然）包括一个预期的语义。
然而，“逻辑”一词的这种技术性用法引出了一个显而易见的问题：
非一阶逻辑的逻辑系统，与直观或哲学意义上的“逻辑”一词究竟有何关联？
下文描述的所有系统皆旨在对某种推理形式进行建模；
我们能否指出，是什么使得它们成为“逻辑的”？

对此没有轻而易举的答案。“逻辑”一词在不同语境中有不同用法，
且与“真理”概念一样，人们已从无数哲学立场对它进行过分析。
例如，人们可能认为逻辑推理的目标在于确定哪些语句必然为真、先验为真、
其真值独立于非逻辑词项的解释、因形式而为真，抑或因语言约定而为真；
而上述每一种观念都需要大量的澄清。
即使我们仅关注数学中所使用的那种逻辑，对其范围也鲜有共识。
例如，在《数学原理》中，罗素与 怀特海
试图在由 弗雷格开创的 {\em 逻辑主义} 传统基础上，
从逻辑出发发展数学。他们的逻辑系统是一种类似于下文所述的高类型逻辑。
最终，他们被迫引入了一些按大多数标准看似乎并非纯逻辑的公理
（特别是无穷公理与可化归公理），
但仍可能有人主张，某些形式的高阶推理应当被承认为逻辑的。
相反，蒯因的本体论并不承认“命题”是合法的论说对象，
他认为二阶与高阶逻辑实质上是披着羊皮的集合论；
换言之，涉及谓词量化的系统并非纯逻辑的。

目前，最好将这些哲学问题留待日后再议，只需将下述系统
视为各类推理（无论逻辑与否）的形式化理想化模型。

\end{document}